\scanpage{49}
\mspage{3}
% archive foliation: 48
les ``diagonales'' $\Delta_1,\Delta_2,\Delta_3$ de $\Delta$ sont déterminées \inserted{(en tant que parties à deux éléments de $\Delta$)} \deleted{par la condition que leur stabilisateur soit} comme les classes pour la relation d'équivalence $R$ définie par l'application $s\longmapsto\Gamma_s$ ($s$ et $s'$ sont dans une $=$ classe ssi ils ont le $=$ groupe d'isotropie, qui est \inserted{cyclique} d'ordre $4$), et la structure polyédrale est déterminée en termes de cette décomposition de $\Delta$, les arêtes étant les \inserted{12} parties à deux éléments distinctes des $\Delta_i$, et les \inserted{8} faces étant les parties à trois éléments \uline{sections} de $\Delta\longrightarrow\Delta/R$.

Donc la structure octaédrale sur $\Delta$ revient à la structure définie par $R$, relation d'équivalence de degré $2$, ou encore par celle du groupe $\Gamma\ (\simeq\mathfrak S_4)$ opérant sur $\Delta$.

Cette procédure combinatoire montre que la structure octaédrale \unsure{annoncée} sur $\Delta\subset X$ est nécessairement unique (puisqu'on connaît son schéma en groupes d'autom.), ce qui montre que la question d'existence est locale (ét.), donc \unsure{OPS}
\[
X=\mathbb P^1_S,\qquad\Delta=\{0,\infty,1,-1,i,-i\}.
\]
Le \deleted{stabilisateur de $0,\infty$} \inserted{sous-groupe de} $\unsure{\Gamma}$ qui fixe $0$ et $\infty$ est isomorphe visiblement à
\[
\mu_4\deleted{\unsure{\simeq}}\simeq(\mathbb Z/4\mathbb Z)_S
\]
(qui contient $i$), \deleted{et par} et de $=$ (par raison d'homogénéité) pour $(1,-1)$, $(i,-i)$. Par contre, le
